\chapter{Knot invariants and the Kontsevich integral}
\label{ch:kontsevich-integral}

\index{Kontsevich integral|textbf}
\index{Vassiliev invariants|textbf}

Configuration-space integrals compute knot invariants only after a
real-locus passage.  The logarithmic bar form
$\eta_{ij}=d\log(z_i-z_j)$ lives on complex configuration space.  For
an embedded real circle $\gamma:S^1\hookrightarrow X$ in a complex
curve, the comparison map is the pullback
\[
\rho_\gamma:\Omega^\bullet_{\log}(\overline{C}_m(X))
\longrightarrow
\Omega^\bullet(\overline{C}_m(S^1)),
\qquad
\rho_\gamma(\omega)=(\gamma^{\times m})^*\omega ,
\]
followed by angular projection and Kontsevich's chord-diagram
quotient.  The ambient bar complex does not by itself produce a
universal finite-type invariant; it produces logarithmic forms and
OPE contractions whose real-locus restriction gives the integrand of
the Kontsevich integral.

The common geometry is Fulton--MacPherson compactification with
Stokes boundary terms.  The 4T relation governing chord diagrams is
the image, under OPE contraction and real-locus restriction, of the
Arnold--Orlik--Solomon relation for logarithmic forms.  A metrized
Lie algebra supplies the Casimir contractions; a representation
supplies the trace around the oriented circle.  This is the exact
surface on which affine Kac--Moody OPE data become Lie algebra weight
systems.

The comparison has three pieces: the real-locus restriction
$\rho_\gamma$, the normalization
$\omega_{ij}^{K}=(2\pi)^{-1}d\arg(z_i-z_j)$, and the
Fulton--MacPherson boundary correction measuring the exact radial
part of $d\log(z_i-z_j)$ (Proposition~\ref{prop:propagator-restriction}).

\section{Vassiliev invariants: mathematical setup}
\label{sec:vassiliev-setup}
\index{Vassiliev invariants!mathematical setup}

\begin{definition}[Vassiliev invariant]
\label{def:vassiliev-invariant}
\index{Vassiliev invariant|textbf}
A knot invariant $v\colon \{\text{oriented knots in } S^3\} \to \bC$
is a \emph{Vassiliev invariant of order~$n$} (or \emph{finite-type
invariant of type~$n$}) if it vanishes on all singular knots with
$n+1$ double points. Equivalently, if $K_0, K_+, K_-$ denote knots
that differ only at a single crossing, the \emph{Vassiliev relation}
$v(K_\times) := v(K_+) - v(K_-)$ extends $v$ to singular knots,
and $v$ has type~$n$ if $v(K)=0$ for every singular knot $K$ with
$n+1$ transverse double points.
\end{definition}

\begin{definition}[Chord diagram]
\label{def:chord-diagram}
\index{chord diagram|textbf}
A \emph{chord diagram of order~$n$} is an oriented circle with
$n$ disjoint pairs of marked points (``chords''), considered up to
orientation-preserving diffeomorphism of the circle. The vector
space $\cD_n$ of chord diagrams is the formal $\bC$-span of
chord diagrams of order~$n$, modulo the four-term
(\emph{4T}) relation.
\end{definition}

\begin{theorem}[Kontsevich integral \cite{Kon94,BarNatan95}]
\label{thm:kontsevich-invariant}
\index{Kontsevich integral!universality}
Let $K:S^1\hookrightarrow \bR^3$ be an oriented framed knot and choose
a height function and complex projection for Kontsevich's
configuration integrals.  After the standard framing normalization
and the Drinfeld associator determined by the regularization, the series
\[
Z_K(K) \;=\;
\sum_{n=0}^\infty \;
\frac{1}{(2\pi i)^n}
\int_{\substack{t_{\min} < t_1 < \cdots < t_n < t_{\max} \\
 z_i = K(t_i),\; z'_i = K(t'_i)}}
\bigwedge_{j=1}^n
\widetilde\omega^K_{j}
\;\cdot\; D_\Gamma
\]
defines a group-like element in the completed chord-diagram algebra
and is universal for Vassiliev invariants: every finite-type invariant
of type~$n$ factors through the degree-$\leq n$ truncation of $Z_K$.
Here $\widetilde\omega^K_{j}=d\log(z_j-z'_j)$ is the unnormalized
logarithmic form; the normalized angular propagator is
$(2\pi)^{-1}d\arg(z_j-z'_j)$.  The diagram $D_\Gamma$ is determined
by the pairing $(t_j,t'_j)$.

Equivalently, the coefficient of a chord diagram is obtained by the
chain
\[
\Omega^\bullet_{\log}(\overline{C}_{2n}(\bC))
\xrightarrow{\ \rho_K\ }
\Omega^\bullet(\overline{C}_{2n}(S^1))
\xrightarrow{\ \int_{\Delta_\Gamma(K)}\ }
\bC\,D_\Gamma ,
\]
where $\rho_K$ is pullback along the knot parameters and
$\Delta_\Gamma(K)$ is the compactified simplex of the ordered chord
endpoints.  This restriction-to-$S^1$ step is part of the theorem.
\end{theorem}

\begin{definition}[Weight system]
\label{def:weight-system}
\index{weight system|textbf}
A \emph{weight system of degree~$n$} is a linear functional
$W\colon \cD_n \to \bC$ satisfying the \emph{4T relation}.  Let
$(\fg,\kappa)$ be a finite-dimensional Lie algebra with
non-degenerate invariant form, let $\Omega=\sum_a x_a\otimes x^a$
be the corresponding Casimir, and let $V$ be a finite-dimensional
$\fg$-module.  For a chord diagram $\Gamma$ with endpoints
$1,\ldots,2n$ in cyclic order, the Lie algebra weight system is
\[
W_{\fg,V}(\Gamma)
\;=\;
\sum_{\{a_e,b_e\}}
\left(\prod_{e\in E(\Gamma)}\kappa^{a_e b_e}\right)
\operatorname{Tr}_{V}\!\left(
X_{c_1}\cdots X_{c_{2n}}
\right),
\]
where $c_r=a_e$ or $b_e$ according as the $r$-th endpoint is the
first or second endpoint of the chord $e$, and $X_c$ denotes the
action of $x_c$ on $V$.  For Jacobi diagrams with internal trivalent
vertices, the same rule inserts the structure tensor $f^a_{bc}$ at
each trivalent vertex.  Invariance of $\kappa$ and the Jacobi
identity give the 4T relation.
\end{definition}

\section{Weight systems from the bar complex}
\label{sec:weight-systems-bar}
\index{weight system!from bar complex}

The bar differential extracts the Casimir and Jacobi tensors that
define a weight system.  Knot invariants appear only after pulling the
forms to the knot circle and applying the Kontsevich integral.

\begin{theorem}[Real-locus bar weight systems]
\label{thm:bar-weight-systems}
Let $\cA=V_k(\fg)$ be the affine Kac--Moody vertex algebra at
non-critical level $k\neq -h^\vee$, normalized so that
\[
J^a(z)J^b(w)
\sim
\frac{k\kappa^{ab}}{(z-w)^2}
\;+\;
\frac{f^{ab}_{c}J^c(w)}{z-w}.
\]
Let $\gamma:S^1\hookrightarrow \bP^1$ be a smooth embedded real
circle contained in an affine chart and let
\[
\rho_\gamma:\Omega^\bullet_{\log}(\overline{C}_{2n}(\bP^1))
\longrightarrow
\Omega^\bullet(\overline{C}_{2n}(S^1))
\]
be the real-locus restriction.  For each degree-$n$ chord diagram
$\Gamma$, the restricted angular form
\[
\omega^K_\Gamma
\;=\;
\bigwedge_{(i,j)\in E(\Gamma)}
\frac{1}{2\pi}\,d\arg(\gamma_i-\gamma_j)
\]
together with the affine OPE contractions defines the Lie algebra
weight system $W_{\fg,V}$.  With $\kappa$ fixed as above, every chord
contributes the factor $k$, so
\[
W_{k\kappa,V}(\Gamma)=k^n W_{\kappa,V}(\Gamma).
\]
The construction factors as
\[
\barB^{(0)}(V_k(\fg))
\longrightarrow
\Omega^\bullet_{\log}(\overline{C}_{2n}(\bP^1))\otimes U(\fg)^{\otimes 2n}
\xrightarrow{\rho_\gamma}
\Omega^\bullet(\overline{C}_{2n}(S^1))\otimes U(\fg)^{\otimes 2n}
\xrightarrow{\operatorname{Tr}_V}
\cD_n^\vee .
\]
It supplies weight systems.  Their composition with
Theorem~\ref{thm:kontsevich-invariant} supplies finite-type knot
invariants.
\end{theorem}

\begin{proof}
The genus-$0$ bar differential uses collision residues of
$\eta_{ij}=d\log(z_i-z_j)$ on
$\overline{C}_{2n}(\bP^1)$ (Definition~\ref{def:bar-differential-complete}).
The logarithmic kernel lowers the OPE pole order by one
(Remark~\ref{rem:propagator-weight-universality}).  Hence the
double-pole term $k\kappa^{ab}(z-w)^{-2}$ contributes the Casimir
edge $k\Omega$, while the simple-pole term contributes the bracket
tensor $f^{ab}_{c}$.  Pullback by $\rho_\gamma$ does not alter these
Lie tensors; it changes only the configuration-space form, replacing
the ambient logarithmic form by its angular Kontsevich component as
computed in Proposition~\ref{prop:propagator-restriction}.

The 4T relation is the image of the Arnold relation
\[
\eta_{ij}\wedge\eta_{jk}
+\eta_{jk}\wedge\eta_{ki}
+\eta_{ki}\wedge\eta_{ij}=0
\]
under the same restriction and OPE contraction.  On the Lie side this
is exactly invariance of $\kappa$ together with the Jacobi identity;
on the configuration-space side it is the codimension-$2$ Stokes
cancellation on the Fulton--MacPherson boundary.  Thus the resulting
linear functional on chord diagrams is a weight system.  Universality
is not asserted at the bar-complex stage; it enters through
Kontsevich's universal integral, whose coefficients are obtained only
after the $\overline{C}_{2n}(\bP^1)\to\overline{C}_{2n}(S^1)$
restriction and chord projection.
\end{proof}

This theorem refines Proposition~\ref{prop:vassiliev-genus0} by
making the real-locus and chord-diagram maps explicit.

\section{The Kontsevich propagator from holomorphic restriction}
\label{sec:kontsevich-propagator}
\index{Kontsevich integral!propagator restriction}

The holomorphic propagator restricts to the Kontsevich propagator
after angular projection.  Its radial part is exact on the open
configuration chamber and contributes only through
Fulton--MacPherson boundary transgression.

\begin{proposition}[Angular restriction and boundary transgression]
\label{prop:propagator-restriction}
Let $S^1 = \{|z| = 1\} \subset \bC$ be the unit circle, parametrized
by $z = e^{i\theta}$. The FM propagator $\eta_{ij} = d\log(z_i-z_j)$
restricts on each ordered chamber of $C_n(S^1)$ as
\[
\eta_{ij}\big|_{S^1}
\;=\;
\frac12\cot\!\left(\frac{\theta_i-\theta_j}{2}\right)
(d\theta_i-d\theta_j)
\;+\;
\frac{i}{2}(d\theta_i+d\theta_j).
\]
Equivalently,
\[
\eta_{ij}\big|_{S^1}
\;=\;
d\log|e^{i\theta_i}-e^{i\theta_j}|
\;+\;
i\,d\arg(e^{i\theta_i}-e^{i\theta_j}).
\]
Thus the normalized Kontsevich angular propagator is
\[
\omega^K_{ij}
\;=\;
\frac{1}{2\pi}\,d\arg(e^{i\theta_i}-e^{i\theta_j})
\;=\;
\frac{1}{2\pi}\,\Im\big(\eta_{ij}\big|_{S^1}\big),
\]
and the normalized holomorphic form decomposes as
\[
\frac{1}{2\pi i}\,\eta_{ij}\big|_{S^1}
\;=\;
\omega^K_{ij}
\;+\;
\frac{1}{2\pi i}\,d\log|e^{i\theta_i}-e^{i\theta_j}|.
\]
For a graph $\Gamma$ and a compactified configuration chain
$F_\Gamma\subset\overline{C}_{V(\Gamma)}(S^1)$,
\[
\int_{F_\Gamma}\bigwedge_{e\in E(\Gamma)}
\frac{1}{2\pi i}\rho_{S^1}(\eta_e)
\;-\;
\int_{F_\Gamma}\bigwedge_{e\in E(\Gamma)}\omega^K_e
\;=\;
\int_{\partial F_\Gamma} T_\Gamma ,
\]
where $T_\Gamma$ is the transgression form obtained from the path
$\omega_e(s)=\omega^K_e+s(2\pi i)^{-1}d\log|z_{h(e)}-z_{t(e)}|$:
\[
T_\Gamma
\;=\;
\int_0^1\sum_{e\in E(\Gamma)}
(-1)^{\epsilon(e)}
\frac{\log|z_{h(e)}-z_{t(e)}|}{2\pi i}
\bigwedge_{e'\neq e}\omega_{e'}(s)\,ds .
\]
Here $\epsilon(e)$ is the Koszul sign from moving the displayed
$0$-form past the preceding edge forms.  Equality of holomorphic and
Kontsevich graph weights is therefore equality after angular
projection plus the vanishing, cancellation, or framing
renormalization of the boundary terms $\int_{\partial F_\Gamma}T_\Gamma$.
\end{proposition}

\begin{proof}
On a chamber with fixed cyclic order,
\[
e^{i\theta_i}-e^{i\theta_j}
=2i\,e^{i(\theta_i+\theta_j)/2}
\sin\!\left(\frac{\theta_i-\theta_j}{2}\right).
\]
Taking $d\log$ gives the displayed formula.  The second displayed
identity is the polar decomposition
$d\log u=d\log|u|+i\,d\arg(u)$ for $u\in\bC^\times$.
The normalized angular form is therefore
$(2\pi)^{-1}d\arg(u)$, while the difference from
$(2\pi i)^{-1}d\log(u)$ is exact on the chamber.

For a product over the edges of a graph, set
$\lambda_e=(2\pi i)^{-1}d\log|z_{h(e)}-z_{t(e)}|$ and
$\omega_e(s)=\omega^K_e+s\lambda_e$.  Since both $\omega^K_e$ and
$\lambda_e$ are closed and $\lambda_e=d((2\pi i)^{-1}
\log|z_{h(e)}-z_{t(e)}|)$,
\[
\frac{d}{ds}\bigwedge_e\omega_e(s)=dT_\Gamma(s).
\]
Integrating in $s$ and applying Stokes on the manifold with corners
$F_\Gamma$ gives the boundary formula.  The Kontsevich framing
correction is precisely the choice of normalization that absorbs the
surviving one-chord boundary term.
\end{proof}

\begin{remark}[Angular and radial parts]\label{rem:imaginary-part}
The angular part $(2\pi)^{-1}d\arg(z_i-z_j)$ is the topological
Kontsevich propagator.  The radial part
$(2\pi i)^{-1}d\log|z_i-z_j|$ is exact on each chamber but not
invisible on the compactification: Stokes moves it to collision and
endpoint faces.  The framing correction is the normalization of these
boundary contributions.  The bar complex produces the full
holomorphic form; the Kontsevich invariant uses its angular
projection together with this boundary normalization.
\end{remark}


\section{The Drinfeld associator from affine bar monodromy}
\label{sec:drinfeld-associator}
\index{Drinfeld associator|textbf}
\index{KZ equation!from bar complex}

The affine genus-$0$ bar comparison carries the KZ connection.  With
tangential basepoints and a logarithm branch fixed, its regularized
monodromy is the KZ Drinfeld associator.

\begin{proposition}[KZ connection from affine bar residues]
\label{prop:kz-from-bar}
\index{KZ equation!from bar differential}
For $\cA=V_k(\fg)$ at non-critical level $k\neq -h^\vee$, the
genus-$0$ affine bar residue comparison on
$C_n(\bC)\subset\overline{C}_n(\bP^1)$ gives the flat connection on
the trivial bundle with fiber $U(\fg)^{\otimes n}$:
\begin{equation}\label{eq:kz-connection}
\nabla^{\mathrm{KZ}}
\;=\;
d \;-\; \frac{1}{k+h^\vee}
\sum_{1 \leq i < j \leq n}
\Omega_{ij}\,d\log(z_i-z_j),
\end{equation}
where $\Omega_{ij} = \sum_a J^a_i J^a_j$ is the Casimir element
acting on the $i$-th and $j$-th tensor factors. This is the
Knizhnik--Zamolodchikov connection.
\end{proposition}

\begin{proof}
The bar differential for $V_k(\fg)$ at genus~$0$
(Definition~\ref{def:bar-differential-complete}) is
$d = \sum_{i<j} \Res_{D_{ij}} (\eta_{ij} \wedge -)$,
where $\eta_{ij} = d\log(z_i-z_j)$ and $\Res_{D_{ij}}$
extracts the Poincar\'e residue along the collision divisor.
The $d\log$ kernel absorbs one pole order from the OPE
(Remark~\ref{rem:propagator-weight-universality}),
so the OPE $J^a(z_i) J^b(z_j) \sim k\kappa^{ab}/(z_i-z_j)^2
+ f^{ab}_c J^c/(z_i-z_j)$ contributes a simple-pole term
$k\kappa^{ab}/(z_i-z_j)$ (Killing contraction) and a regular
term $f^{ab}_c J^c$ (adjoint action) to the collision residue.
Contracting with the Casimir gives
$\Omega_{ij}\,d\log(z_i-z_j)$. The factor $1/(k+h^\vee)$ arises from
the Sugawara normalization: the OPE is formulated at level~$k$,
and the KZ connection uses the normalized Casimir
$\Omega/(k+h^\vee)$. (At the critical level $k = -h^\vee$,
the Sugawara construction is undefined and the KZ connection
degenerates; the correct replacement is the Feigin--Frenkel
connection on opers.)

Flatness $(\nabla^{\mathrm{KZ}})^2 = 0$ follows from $d^2 = 0$
on the genus-$0$ bar complex (the bar differential restricted to
the open stratum $C_n(\bC)$ has no curvature since $\kappa$ arises
only at genus~$\geq 1$). Explicitly, the flatness condition is
the infinitesimal braid relation
$[\Omega_{ij}, \Omega_{ik} + \Omega_{jk}] = 0$ for $i,j,k$ distinct,
which follows from the Jacobi identity for~$\fg$.
\end{proof}

\begin{theorem}[KZ associator on the bar comparison surface]
\label{thm:drinfeld-associator-bar}
\index{Drinfeld associator!from bar monodromy}
The monodromy of the KZ connection~\eqref{eq:kz-connection}
on $C_3(\bC)$ around the boundary divisors of
$\overline{C}_3(\bP^1)$ produces the Drinfeld associator
$\Phi_{\mathrm{KZ}}\in U(\fg)^{\otimes 3}[[\hbar]]$ with
\[
h=(k+h^\vee)^{-1},
\qquad
\hbar=2\pi i\,h .
\]
The statement uses the standard tangential basepoints at $0$ and $1$;
changing the associator choice translates
$\Phi_{\mathrm{KZ}}$ by the Drinfeld $\mathrm{GRT}$-torsor.  Under
the affine bar comparison:
\begin{enumerate}[label=\textup{(\roman*)}]
\item The parallel transport of the bar complex
 $\barB(V_k(\fg))$ around the loop
 $\{z_1 = 0,\, z_2 = 1,\, z_3 = t\}$ for $t \in (0,1)$
 with $t \to 0$ and $t \to 1$ as the two boundary degenerations
 gives the regularized holonomy
 $\Phi_{\mathrm{KZ}}(t_{12}, t_{23})$
 where $t_{ij}=h\Omega_{ij}$.
\item The pentagon relation for the associator
 is a consequence of the coherence of the bar complex
 on $C_4(\bC)$: the five boundary strata of
 $\overline{\mathcal{M}}_{0,5} \cong \operatorname{Bl}_4(\bP^2)$
 give five ways to degenerate four-point configurations,
 and the associator intertwines them.
\item The hexagon relation comes from the braid monodromy
 on $C_3(\bC)$: the half-twist $\sigma_1$ exchanges
 $z_1$ and $z_2$, and
 the corresponding braiding is
 $R_{12}=\exp(\pi i\,h\Omega_{12})$.
\end{enumerate}
\end{theorem}

\begin{proof}
The Drinfeld--Kohno theorem \cite{Drinfeld90,Kohno87} establishes
that the regularized monodromy of the KZ connection produces the
associator.  Proposition~\ref{prop:kz-from-bar} identifies this
connection form with the genus-$0$ affine bar residue comparison.

Part (i): the bar differential on $C_3(\bC)$ is
$d = t_{12}\, \eta_{12} + t_{13}\, \eta_{13} + t_{23}\, \eta_{23}$
where $\eta_{ij} = d\log(z_i - z_j)$ and
$t_{ij}=h\Omega_{ij}$. The parallel transport from
$t = \varepsilon$ to $t = 1 - \varepsilon$ is
$\Phi_{\mathrm{KZ}}(t_{12}, t_{23})
= 1 + \sum_{n \geq 1} \int_{0<s_1<\cdots<s_n<1}
\omega(s_1) \cdots \omega(s_n)$
where $\omega(s) = t_{12}\, ds/(s) + t_{23}\, ds/(s-1)$, regularized
at $s = 0$ and $s = 1$. This is the iterated integral formula for
the Drinfeld associator.

Part (ii): on $C_4(\bC)$, the bar complex structure requires
coherence of all four-point degenerations. The boundary of
$\overline{\mathcal{M}}_{0,5}$ has five associahedral components,
and the parallel transports around these components compose to give
Drinfeld's pentagon
\[
\Phi^{1,2,34}\Phi^{12,3,4}
\;=\;
\Phi^{2,3,4}\Phi^{1,23,4}\Phi^{1,2,3}.
\]

Part (iii): the half-twist $\sigma_1$ on $C_3(\bC)$ gives
the braid generator, and the monodromy of the KZ connection
under this half-twist is $R_{12}=\exp(\pi i\,h\Omega_{12})$
(the universal $R$-matrix at level~$k$). The hexagon relation
expresses compatibility of the associator with the braiding.
\end{proof}

\begin{remark}[Bar-cobar as quasi-triangular structure]
\label{rem:bar-quasi-triangular}
\index{quasi-triangular!from bar complex}
Theorem~\ref{thm:drinfeld-associator-bar} supplies the
quasi-triangular quasi-Hopf data on the affine Kac--Moody comparison
surface: the associator from three-point KZ monodromy, the
$R$-matrix from two-point braid monodromy, and the pentagon/hexagon
from the boundary stratification of $C_4(\bC)$ and $C_3(\bC)$.  The
input is the fixed KZ associator choice above.  The
Reshetikhin--Turaev package is obtained after applying the
Drinfeld--Kohno theorem and the chosen associator; it is not an
unqualified equality between an arbitrary bar complex and an arbitrary
quantum group presentation.
\end{remark}


\section{Higher genus: graph complex and real-loop comparison}
\label{sec:bar-to-loop}
\index{Kontsevich integral!loop expansion}

Positive genus separates into an algebraic graph-complex theorem and
a real-propagator comparison class.  The first is a statement about
the modular Feynman transform.  The second is a statement about
pulling prime-form propagators to a real cycle and comparing the
result with the topological propagator of perturbative
Chern--Simons theory.

\begin{proposition}[Feynman transform and algebraic graph-complex
refinement]
\label{prop:feynman-graph-complex}
\index{Feynman transform!graph complex}
The Feynman transform of the modular operad $\mathrm{Com}$ gives the
stable graph complex of Getzler--Kapranov.  For a Koszul chiral
algebra $\cA$, the full bar complex is the same stable-graph
construction with vertices decorated by the OPE tensors of~$\cA$:
\begin{enumerate}[label=\textup{(\roman*)}]
\item Getzler--Kapranov~\cite[Theorem~5.4]{GeK98} identify
 $\mathrm{FT}(\mathrm{Com})$ with the modular stable-graph complex.
\item The identification
 $\barB^{\mathrm{full}}(\cA)
 \simeq \mathrm{FT}_{\mathcal{M}\mathrm{od}}(\cA)$
 (Theorem~\ref{thm:prism-higher-genus}) decorates the same graph
 complex by chiral OPE data.
\item The passage from this algebraic graph complex to the loop
 expansion of the Kontsevich integral is the real-locus restriction
 and boundary-normalization problem of
 Definition~\ref{def:real-loop-obstruction}; it is not a consequence
 of the Feynman-transform isomorphism alone.
\end{enumerate}
\end{proposition}

\begin{proof}
Part (i) is Getzler--Kapranov \cite[Theorem~5.4]{GeK98}.
Part~(ii): Theorem~\ref{thm:prism-higher-genus} identifies the
genus-$g$ bar complex with the genus-$g$ component of the
Feynman transform $\mathrm{FT}_{\mathcal{M}\mathrm{od}}(\cA)$.
At $\cA = \mathrm{Com}$ (the commutative operad, viewed as a
trivial chiral algebra), the OPE data is trivial and
$\mathrm{FT}_{\mathcal{M}\mathrm{od}}(\mathrm{Com}) =
\mathrm{FT}(\mathrm{Com}) = \mathrm{GC}$.
Part~(iii) follows from functoriality of differential forms: a graph
weight on $\overline{C}_{m}(\Sigma_g)$ becomes a knot-integral weight
only after pullback to $\overline{C}_m(S^1)$ and quotient by the
chord-diagram relation.  Proposition~\ref{prop:propagator-restriction}
shows already in genus~$0$ that the radial exact term contributes
boundary transgressions; positive genus has the same obstruction with
the prime-form propagator replacing $d\log(z-w)$.
\end{proof}

\begin{definition}[Real-loop comparison obstruction]
\label{def:real-loop-obstruction}
\label{rem:holomorphic-to-topological}
\index{Kontsevich integral!holomorphic vs topological}
Let $\Sigma_g$ be a compact Riemann surface, let
\(\gamma:S^1\hookrightarrow\Sigma_g\) be an embedded real cycle, and
let
\[
\omega^{\mathrm{hol}}_{ij}
\;=\;
d_{z_i}\log E(z_i,z_j)
\]
be the prime-form logarithmic propagator on the complement of the
diagonal.  For a graph $\Gamma$ define
\[
\mathfrak{o}^{g,\gamma}_\Gamma
\;=\;
\left[
\rho_\gamma\!\left(\bigwedge_{e\in E(\Gamma)}
\omega^{\mathrm{hol}}_e\right)
\;-\;
\bigwedge_{e\in E(\Gamma)}\omega^K_e
\;-\;dT^{g,\gamma}_\Gamma
\right]
\in
H^{|E(\Gamma)|}\!\left(
\overline{C}_{V(\Gamma)}(S^1),\partial;\bC\right),
\]
where $T^{g,\gamma}_\Gamma$ is the Fulton--MacPherson transgression
over all collision and endpoint faces, normalized as in
Proposition~\ref{prop:propagator-restriction}.  The positive-genus
bar graph weight for $\Gamma$ agrees with the topological
Kontsevich/Chern--Simons graph weight along $\gamma$ exactly when
\(\mathfrak{o}^{g,\gamma}_\Gamma=0\) and the chosen trivialization of
the transgression matches the Chern--Simons gauge-fixing
normalization of Cattaneo--Mnev~\cite{CattaneoMnev10}.
\end{definition}


\section{Chern--Simons comparison class}
\label{sec:cs-bridge}
\index{Chern--Simons!as bridge to Vassiliev}
Chern--Simons gauge theory supplies the topological configuration-space
model.  The comparison with the affine bar complex is a morphism of
factorization algebras plus a Wilson-line trace comparison, controlled
by an obstruction class.

\begin{definition}[Chern--Simons boundary comparison class]
\label{def:cs-factorization-comparison}
\index{Chern--Simons!factorization homology}
Let $M$ be a compact oriented $3$-manifold with boundary
$\partial M=\Sigma$, let $\mathrm{Obs}^{q}_{\mathrm{CS},k}$ be the
perturbative Chern--Simons factorization algebra at level~$k$ in the
sense of Costello--Gwilliam~\cite{CG17}, and let
\[
\mathcal B_{\Sigma,k}:=\barB(V_k(\fg))_\Sigma
\]
be the affine bar factorization algebra on the boundary Riemann
surface.  A collar trivialization and a boundary polarization define a
local comparison morphism
\[
\beta_{\partial}:
\mathrm{Obs}^{q}_{\mathrm{CS},k}\big|_{\Sigma\times[0,\varepsilon)}
\longrightarrow
\mathcal B_{\Sigma,k}.
\]
The Chern--Simons boundary comparison class is
\[
\mathfrak{o}_{\mathrm{CS}}(M,\Sigma,k)
\in
H^1\!\operatorname{Cone}\!\left(
C^\bullet_{\mathrm{fact}}\!
\left(\mathrm{Obs}^{q}_{\mathrm{CS},k}\big|_{\Sigma\times[0,\varepsilon)}\right)
\xrightarrow{\ \beta_\partial\ }
C^\bullet_{\mathrm{fact}}(\mathcal B_{\Sigma,k})
\right),
\]
the \v{C}ech--BV descent obstruction to extending
$\beta_\partial$ to a global factorization-algebra comparison.  For a
colored knot $(K,V)\subset S^3$, the Wilson trace comparison has the
secondary obstruction
\[
\mathfrak{o}_{K,V}
\;:=\;
\langle W_K(V)\rangle_{\mathrm{CS}}
\;-\;
\chi(\barB(V_k(\fg),V_K))
\in \bC[[\hbar]] .
\]
\end{definition}

\begin{proposition}[Chern--Simons comparison criterion]
\label{prop:cs-comparison-criterion}
With the choices in Definition~\ref{def:cs-factorization-comparison}, the
identity
\[
\langle W_K(V)\rangle_{\mathrm{CS}}
\;=\;
\chi(\barB(V_k(\fg),V_K))
\]
holds on the perturbative comparison surface only after the two
classes
\[
\mathfrak{o}_{\mathrm{CS}}(S^3\setminus N(K),\partial N(K),k),
\qquad
\mathfrak{o}_{K,V}
\]
vanish and the propagator normalization agrees with
Proposition~\ref{prop:propagator-restriction}; here $N(K)$ is a
tubular neighborhood of the Wilson line.  Costello--Gwilliam
\cite{CG17} construct the factorization-algebra side; the
Costello--Francis--Gwilliam comparison with knot polynomials
\cite{CFG25} supplies the Reshetikhin--Turaev trace surface; the
Axelrod--Singer--Kontsevich configuration integral surface is
encoded by restriction of forms
\[
K^*:\Omega^\bullet(\overline{C}_n(S^3))
\longrightarrow
\Omega^\bullet(\overline{C}_n(S^1)),
\]
induced by $K_*:\overline{C}_n(S^1)\to\overline{C}_n(S^3)$ and
described in Remark~\ref{rem:vassiliev}.
\end{proposition}


\section{Kontsevich formality and the \texorpdfstring{$E_2$}{E2} obstruction}
\label{sec:kontsevich-formality-rectification}
\index{Kontsevich formality!$E_2$ obstruction|textbf}
\index{E2-formality@$E_2$-formality!Kontsevich}

\begin{construction}[Kontsevich formality as the genus-$0$ $E_2$ model]
\label{constr:kontsevich-rectification}
\index{formality!Kontsevich|textbf}
\index{polyvector fields!formality}
\index{polydifferential operators!formality}
Kontsevich's formality theorem, the $L_\infty$
quasi-isomorphism
$\Phi \colon T_{\mathrm{poly}} \xrightarrow{\sim}
D_{\mathrm{poly}}$, strictifies the homotopy Gerstenhaber algebra of
polyvector fields to the dg algebra of polydifferential operators.
In the modular Koszul framework this is $E_2$-formality of the
genus-$0$ convolution algebra $\mathfrak{g}^{(0)}_\cA$
(Proposition~\ref{prop:en-formality-mc-truncation}\textup{(ii)}).
\begin{enumerate}[label=\textup{(\roman*)}]
\item \emph{Classical data.}
 The Hochschild cochain complex $C^*(A,A)$ of an associative
 algebra~$A$ carries a $G_\infty$-structure (homotopy
 Gerstenhaber): cup product~$\cup$ and Gerstenhaber
 bracket~$[-,-]$ satisfying Leibniz only up to coherent
 homotopies. The polyvector fields $T_{\mathrm{poly}}(\bR^d)$
 carry the strict Schouten--Nijenhuis bracket and wedge product.
 Kontsevich's maps $U_n \colon T_{\mathrm{poly}}^{\otimes n}
 \to D_{\mathrm{poly}}$ are given by integrals over
 $\overline{C}_n(\mathfrak{H})$ (the Fulton--MacPherson
 compactification of the upper half-plane), using the angular
 propagator $d\arg(z_i-z_j)$ that appears in the real-locus
 comparison of Proposition~\ref{prop:propagator-restriction}.
\item \emph{Bar-complex interpretation.}
 The configuration space
 $\overline{C}_n(\mathfrak{H})$ is the genus-$0$, $n$-point FM
 compactification with boundary. The formality maps $U_n$ are
 fiber integrals of bar-complex forms over the FM boundary
 strata, identical in structure to the genus-$0$ bar differential
 of an affine vertex algebra
 (Theorem~\ref{thm:bar-weight-systems}). The Stokes theorem on
 $\overline{C}_n(\mathfrak{H})$ that proves $\Phi$ is an
 $L_\infty$-morphism is the same codimension-$1$ boundary
 cancellation that gives $d^2 = 0$ for the bar complex.
\item \emph{Modular Koszul placement.}
 The $E_2$-formality hierarchy
 (Proposition~\ref{prop:en-formality-mc-truncation}\textup{(ii)})
 states: $E_2$-formality of the genus-$0$ convolution algebra
 $\mathfrak{g}^{(0)}_\cA$ is equivalent to termination of the
 shadow obstruction tower at degree~$2$. Kontsevich formality is exactly
 this statement for the associative operad, where the genus-$0$
 convolution algebra is the Hochschild cochain complex with its
 $G_\infty$-structure. The higher formality maps $U_n$
 ($n \geq 2$) are the transferred $L_\infty$ brackets
 $\ell_n^{(0),\mathrm{tr}}$ of the bipartite shadow
 (Theorem~\ref{thm:bipartite-linfty-tree}).
\end{enumerate}
\end{construction}

\section{Summary}
\label{sec:kontsevich-summary}

\noindent
\textbf{Comparison surface.}
\begin{enumerate}
\item The affine bar complex produces Lie algebra weight systems
 after the real-locus restriction
 $\overline{C}_{2n}(\bP^1)\to\overline{C}_{2n}(S^1)$
 (Theorem~\ref{thm:bar-weight-systems}).
\item The holomorphic propagator decomposes on $S^1$ into the
 angular Kontsevich propagator plus an exact radial term whose
 contribution is the Fulton--MacPherson boundary transgression
 (Proposition~\ref{prop:propagator-restriction}).
\item The genus-$0$ affine bar residue comparison on $C_n(\bC)$ is
 the KZ connection for $V_k(\fg)$
 (Proposition~\ref{prop:kz-from-bar}).
\item The KZ Drinfeld associator arises from regularized affine bar
 monodromy on $C_3(\bC)$ after a tangential-basepoint and associator
 choice (Theorem~\ref{thm:drinfeld-associator-bar}).
\item The Feynman transform of $\mathrm{Com}$ is the stable graph complex
 (Proposition~\ref{prop:feynman-graph-complex}).
\item Positive-genus comparison with the Kontsevich/Chern--Simons
 loop expansion is controlled by the classes
 $\mathfrak{o}^{g,\gamma}_\Gamma$ and
 $\mathfrak{o}_{\mathrm{CS}}$
 (Definition~\ref{def:real-loop-obstruction} and
 Definition~\ref{def:cs-factorization-comparison}).
\end{enumerate}
